% Source: https://artofproblemsolving.com/community/c3h1389217_tips_on_infinite_square_roots
% Stable identity: aops-topic-1389217-q02
Let $L=\dfrac{1+\sqrt{21}}{2}$, the positive root of $L^2=L+5$.

\textbf{The sequence is increasing and bounded by $L$.} We have $0<b_1=\sqrt5<L$ (since $L>2.79$). If $0<b_n<L$, then $b_{n+1}=\sqrt{5+b_n}<\sqrt{5+L}=L$, and $b_{n+1}>b_n$ because $b_{n+1}^2-b_n^2=5+b_n-b_n^2>0$ for $0<b_n<L$ (the quadratic $5+t-t^2$ is positive between its roots $\dfrac{1-\sqrt{21}}{2}<0$ and $L$).

\textbf{The limit.} A bounded increasing sequence converges; its limit $p$ satisfies $p=\sqrt{5+p}$, so $p^2-p-5=0$ and $p=\dfrac{1\pm\sqrt{21}}{2}$. A square root is never negative, so $p\ge0$ and the minus sign is impossible. Hence
\[
\sqrt{5+\sqrt{5+\sqrt{5+\cdots}}}=\boxed{\dfrac{1+\sqrt{21}}{2}}\approx2.79 .
\]
(The post gives both values $\dfrac{1\pm\sqrt{21}}{2}$; only the positive one is correct.)
